\section*{Chapter \ref{ch:iv:statistics} --- Statistics}

\begin{solution}{iq:iv:statistics:1}
The p-value is the probability, computed assuming the null hypothesis is true, of a statistic as far from the null as the observed
one, or farther. ``The probability that the strategy does not work'' is the probability of the null given the data, which a
p-value does not give: that needs a prior, and a small p-value on a strategy chosen from hundreds is still likely to be
a null. It also says nothing about how large the effect is.
\iqlookfor{the conditional in the right direction, and the prior it leaves out.}
\end{solution}

\begin{solution}{iq:iv:statistics:2}
$0.01/\sqrt{250} \approx 0.063\%$ for the mean daily return. The annual mean is 250 times the daily mean, so its
standard error is $250 \times 0.063\% \approx 15.8\%$, the annual volatility itself. One year of data says almost nothing
about the expected return; it takes decades, which is why researchers estimate risk far better than return.
\iqlookfor{the square-root law and its uncomfortable consequence for expected returns.}
\end{solution}

\begin{solution}{iq:iv:statistics:3}
Deviations are taken from the sample mean, which is fitted to the same data and so is closer to them than the true mean;
the sum of squares is too small on average by one variance, and dividing by $n - 1$ makes the estimator unbiased. It
matters when $n$ is small (a few observations per group, short windows) and not at all for thousands of days, where
dependence and fat tails dominate the error.
\iqlookfor{the degree of freedom used by the mean, and a sense of when it matters.}
\end{solution}

\begin{solution}{iq:iv:statistics:4}
The product of the two slopes is $r^2$: $0.8 \times 0.45 = 0.36$, so $r = 0.6$, positive because both slopes are.
\iqlookfor{$b_{y|x} = r s_y/s_x$ and $b_{x|y} = r s_x/s_y$.}
\end{solution}

\begin{solution}{iq:iv:statistics:5}
About 1.16 (\cref{ex:iv:statistics:sharpe}), so $t \approx 1.73$: not significant at 5\% two-sided, borderline one-sided.
If the strategy was one of several considered, or its returns are autocorrelated, the evidence is weaker still.
\iqlookfor{the Sharpe ratio's standard error and the conclusion that nine months prove little.}
\end{solution}

\begin{solution}{iq:iv:statistics:6}
$1 + 2 \times (-0.8)/1 = -0.6$. You would conclude that $x_1$ predicts negatively, when its true effect is positive
and it merely stands in for its negative correlation with $x_2$. Include $x_2$, or residualise $x_1$ on it before
judging.
\iqlookfor{the omitted-variable formula and the practical fix.}
\end{solution}

\begin{solution}{iq:iv:statistics:7}
The slope is multiplied by $\Var(x)/(\Var(x) + \Var(u)) = \tfrac12$: biased towards zero. Correct it with an estimate of
the noise variance (divide the slope by the reliability), with an instrument (a second noisy measurement of the same
signal), or by averaging repeated measurements; a simulation in the chapter's code shows the halving.
\iqlookfor{attenuation bias, its direction and the standard corrections.}
\end{solution}

\begin{solution}{iq:iv:statistics:8}
$\big((1.645 + 0.842)/0.5\big)^2 \approx 24.7$ years; for a Sharpe ratio of 1, about 6.2 years. Track records cannot
establish modest Sharpe ratios; economic reasoning and out-of-sample evidence from other markets have to do much of the
work.
\iqlookfor{a power calculation in annual Sharpe units, and its implication.}
\end{solution}

\begin{solution}{iq:iv:statistics:9}
Five expected false positives; $1 - 0.95^{100} \approx 0.994$ chance of at least one; Bonferroni tests each at $0.05/100 =
0.0005$ to control the family-wise error at 5\%. With many related signals, controlling the false discovery rate
(Benjamini--Hochberg) is usually the better target (One Quant Book~4, chapter~12).
\iqlookfor{expected false positives, the family-wise probability and the corrections.}
\end{solution}

\begin{solution}{iq:iv:statistics:10}
Over five years a Sharpe ratio estimate has standard deviation about $1/\sqrt5 \approx 0.45$, and the best of 200
worthless variants is expected near $2.75 \times 0.45 \approx 1.23$ (the table of the chapter). The reported 1.9 is
about 1.5 standard deviations of one estimate above that. Report the number of variants, the deflated Sharpe ratio,
and the result on data not used to choose the variant; say that 1.9 selected from 200 is weak evidence.
\iqlookfor{the expected maximum under the null and honest reporting.}
\end{solution}

\begin{solution}{iq:iv:statistics:11}
The variance of the mean of an AR(1) with coefficient $\phi$ is inflated by $(1 + \phi)/(1 - \phi) = 3$ for large $n$, so
the naive standard error is too small by $\sqrt3 \approx 1.73$ (the chapter's simulation of 2\,000 series of 500 days
gives about the same). An i.i.d.\ bootstrap reshuffles days and destroys the dependence, reproducing the naive error;
resample blocks of consecutive days (moving-block or stationary bootstrap) with blocks longer than the correlation length,
or use a HAC standard error.
\iqlookfor{the variance inflation factor and a bootstrap that keeps the dependence.}
\end{solution}

\begin{solution}{iq:iv:statistics:12}
Fixed at 12 tosses: $\P(X \le 3) \approx 0.073$ under a fair coin. Tossing until the third head: the evidence is ``it took
12 or more tosses'', $\P(N \ge 12) = \P(\text{at most 2 heads in 11}) \approx 0.033$. The data are identical and the
likelihoods are proportional, but the p-value sums over outcomes that did not occur, which differ between the two plans.
A frequentist accepts that the design matters; a Bayesian (or anyone following the likelihood principle) finds it
troubling. In practice it is a warning about optional stopping: a test run ``until it looks significant'' needs a
sequential method.
\iqlookfor{computing both p-values and understanding why stopping rules enter them.}
\end{solution}

\begin{solution}{iq:iv:statistics:13}
Close to zero: the other 250 days carry no relation, and the chapter's simulation of such a year gives about 0.01 without
the crash day and 0.21 with it. One point with enormous leverage made the fit. Report both, say that the relation is a
crash co-movement and not a daily one, and consider a robust regression or a separate model for extreme days.
\iqlookfor{leverage of a single observation and a report that separates regimes.}
\end{solution}
